% -*-latex-*-
% Non-technical Summary, for a general audience. RMIT requirement: <200 words
% (counted as printed).

A Vietnamese student opens her essay with her grandfather's proverb,
``water wears down stone'', and argues as many Asian traditions
teach: context first, conclusion arrived at rather than announced. A
human examiner sees the skill. An AI grader trained mostly on Western
text may see only a missing thesis statement and mark her down for
the shape of her reasoning.

AI now grades high-stakes essays. This thesis measures that penalty
and proposes CAMPS: a panel of AI judges, each reading the essay
through a different cultural lens, disagreement flagging it for a
teacher.

Every AI model tested scored native-English essays above upper-level
Southeast Asian ones, the widest gap over a full band. Because no
human had judged the two groups to be of equal quality, this is a
disparity, not a proven cultural bias. The panel narrowed the gap by
about a quarter on the two most disparate models, widened it on one,
and more judges did not help; the disagreement signal proved too
coarse, and too loosely tied to cultural disagreement, to pick out
the essays needing a human. Any AI grader's gap can now be measured
before it is trusted; no single cultural lens should grade alone.
